\SourceLeaf{536}vertical column of each individual table the three lowest digits of the square, and in the interior the four or five higher digits of the square.
\par The table ‘\textit{Indices of Primes in the Higher Number Domain}’ gives the modulus in each case in the top horizontal row, the base used in the second row, the residues in the first vertical column, and the indices in the interior. The stars $*$ denote zero residues.
\par The manuscript of the fragment ‘\textit{Eighth Section. A Fuller Consideration of Certain Investigations Concerning Circle Division}’ seems to date from the reworking of the ‘Analysis residuorum’ into the ‘Disquisitiones arithmeticae’. The letters from \textsc{Gauss} to \textsc{Dirichlet} confirm the view that in higher arithmetic, too, \textsc{Gauss} discovered considerably more than is found in his manuscript remains.
\begin{flushright}\textsc{Schering}.\end{flushright}
\BookRule{.25\textwidth}
